% Contribution for the ProveIt continuation. No document preamble.
% Requires amsmath, amssymb, amsthm, booktabs, hyperref.
% References supplied in cm_norms_references.tex.

\section{Exact CM row products from class polynomials}
\label{cmnorm:sec:main}

The manuscript's complex-multiplication chapter records striking products
of weight-four polygamma row sums at discriminants $-15,-20,-23,-39,-47$.
Its displayed numerical residuals establish useful evidence, but do not
explain the prime factors or provide proofs of the gamma products.
The following argument supplies that missing step. It also evaluates
the corresponding products at every even weight by a finite computation
with a Hilbert class polynomial.

The analytic ingredients are classical: the Kronecker limit formula,
the Lerch--Chowla--Selberg formula, and the algebra of modular forms.
The contribution here is an explicit, uniformly normalized consequence
for the manuscript's row sums, accompanied by certified class-polynomial
data and new displayed weight-six evaluations. We do not claim that the
underlying CM period theory is new.

\subsection{Conventions and the main theorem}

Let $D=-d<-4$ be a fundamental discriminant, let
$\chi_D(r)=(D/r)$ be its primitive quadratic Dirichlet character,
and let $h=h(D)$. Take the standard reduced primitive positive definite
forms
\[
 \mathcal Q_D=\{(a,b,c):b^2-4ac=D,\ |b|\leq a\leq c,\
 b\geq0\text{ if }|b|=a\text{ or }a=c\}.
\]
There is one representative for each proper ideal class. Write
\begin{align}
 \tau_Q&=\frac{-b+i\sqrt d}{2a},&
 A_D&=\prod_{Q\in\mathcal Q_D}a,&
 P_D&=\prod_{\substack{1\leq r<d\\\chi_D(r)=1}}
       \Gamma(r/d),\label{cmnorm:eq:data}\\
 H_D(X)&=\prod_{Q\in\mathcal Q_D}(X-j(\tau_Q)),&
 C_d&=\Phi_d(1),&\varphi_d&=\varphi(d).\notag
\end{align}
Thus $H_D\in\mathbb Z[X]$ is monic; $C_d=p$ if $d$ is a positive
power of a prime $p$, and $C_d=1$ otherwise. All fractional powers of
positive real numbers below denote their positive values.

For even $w\geq4$, use exactly the lattice normalization of the manuscript:
\begin{equation}
 G_w(\tau)=\sum_{(m,n)\neq(0,0)}(m+n\tau)^{-w}
 =2\zeta(w)E_w(\tau).
 \label{cmnorm:eq:G}
\end{equation}
The modular discriminant is $\Delta(\tau)=\eta(\tau)^{24}$;
in particular $j=E_4^3/\Delta$ and $j-1728=E_6^2/\Delta$.

\begin{lemma}[The polynomial attached to an even weight]
\label{cmnorm:lem:polynomial}
For every even integer $w\geq4$, there is a unique monic polynomial
$F_w\in\mathbb Q[X]$ of degree $w$ such that
\begin{equation}
 \frac{E_w(\tau)^{12}}{\Delta(\tau)^w}=F_w(j(\tau)).
 \label{cmnorm:eq:Fw}
\end{equation}
It is effectively computable using exact rational Fourier coefficients.
\end{lemma}
\begin{proof}
The left side is a weight-zero modular function, holomorphic on the
upper half-plane since $\Delta$ has no zeros there, with a pole of order
$w$ at the cusp and no other poles. The field of modular functions is
$\mathbb C(j)$, and a modular function with no finite poles is a
polynomial in $j$. Since $E_w=1+O(q)$, $\Delta=q+O(q^2)$, and
$j=q^{-1}+744+O(q)$, this polynomial is monic of degree $w$.
Successively removing its Laurent coefficients at $q^{-w},\ldots,q^0$
computes all coefficients. The input coefficients are rational, hence so
are the polynomial coefficients. Equivalently, one can work in the exact
graded algebra $\mathbb Q[E_4,E_6]$.
\end{proof}

Set
\[
 \mathcal R_{D,w}=\operatorname{Res}_X(H_D(X),F_w(X)).
\]
Our resultant convention is
$\operatorname{Res}(H_D,F_w)=\prod_QF_w(j(\tau_Q))$, since $H_D$
is monic. No unidentified sign or leading coefficient is implicit.

\subsection{An exact generator for the modular weight polynomial}
\label{newmodular:sec:generator}

Normalize the Eisenstein series to have constant coefficient one:
\[
 E_w(\tau)=1-\frac{2w}{B_w}\sum_{m\ge1}\sigma_{w-1}(m)q^m,
 \qquad q=e^{2\pi i\tau},
\]
and set
\[
 \Delta=q\prod_{m\ge1}(1-q^m)^{24}
       =\frac{E_4^3-E_6^2}{1728},\qquad j=\frac{E_4^3}{\Delta}.
\]
For each even $w\ge4$, let $a\in\{0,1,2\}$, $b\in\{0,1\}$, and
$c\ge0$ be the unique integers satisfying
\begin{equation}\label{newmodular:eq:residue}
 w=4a+6b+12c.
\end{equation}
The excluded weight two is precisely the exceptional small weight in
this residue prescription.

The level-one graded-ring structure \cite{cmnorm:Stein} gives a unique monic polynomial
$P_w\in\mathbb Q[X]$ of degree $c$ such that
\begin{equation}\label{newmodular:eq:reduced}
 E_w=E_4^aE_6^b\Delta^cP_w(j).
\end{equation}
This can be computed using only rational arithmetic.  Write
$P_w(X)=\sum_{r=0}^c p_rX^r$ and expand the basis
\[
 E_4^{a+3r}E_6^b\Delta^{c-r},\qquad 0\le r\le c.
\]
The basis element indexed by $r$ has first term $q^{c-r}$ with
coefficient one.  Successive coefficient comparison at
$q^0,q^1,\ldots,q^c$ therefore determines $p_c,p_{c-1},\ldots,p_0$
by unit triangular elimination.  Since $E_w$ has constant term one,
$p_c=1$.

Consequently the polynomial in the CM norm formula is explicitly
\begin{equation}\label{newmodular:eq:F-factor}
 \boxed{\displaystyle
 F_w(X)=X^{4a}(X-1728)^{6b}P_w(X)^{12}.}
\end{equation}
Indeed, raising \eqref{newmodular:eq:reduced} to the twelfth power and
using $E_4^3=j\Delta$, $E_6^2=(j-1728)\Delta$, and
\eqref{newmodular:eq:residue} gives
$E_w^{12}/\Delta^w=F_w(j)$.  The degree of $F_w$ is exactly $w$
and its leading coefficient is one.

\paragraph{Finite certificates.}
The companion generator computes $\Delta$ independently from its Euler
product and verifies its normalization against $E_4^3-E_6^2$.
For every even $w$ from four through twenty-four, it verifies exactly
the coefficients of $q^0,\ldots,q^w$ in
\begin{equation}\label{newmodular:eq:cleared}
 E_w^{12}=\sum_{r=0}^{w}[X^r]F_w(X)\,
                     E_4^{3r}\Delta^{w-r}.
\end{equation}
Both sides are holomorphic modular forms of weight $12w$ on
$\mathrm{SL}_2(\mathbb Z)$.  The level-one Sturm bound is $w$, so
these finitely many rational equalities certify the identity.  In
characteristic zero this bound follows directly from the valence
formula \cite[Theorem 2.11]{cmnorm:Stein}: a nonzero weight-$12w$ holomorphic modular form has order
at infinity at most $w$.

The JSON artifact stores the exact rational coefficients, the compact
factorization, the degree and monicity checks, and the finite
coefficient certificates.  It contains the full range, including
\begin{align}
 F_{14}(X)&=X^8(X-1728)^6,\label{newmodular:eq:F14}\\
 F_{16}(X)&=X^4\left(X-\frac{3456000}{3617}\right)^{12},
 \label{newmodular:eq:F16}\\
 F_{18}(X)&=(X-1728)^6\left(X-\frac{9504000}{43867}\right)^{12},\\
 F_{20}(X)&=X^8\left(X-\frac{209520000}{174611}\right)^{12},\\
 F_{22}(X)&=X^4(X-1728)^6
                  \left(X-\frac{35424000}{77683}\right)^{12},\\
 F_{24}(X)&=\left(X^2-\frac{340364160000}{236364091}X
                         +\frac{30710845440000}{236364091}\right)^{12}.
\end{align}
For example,
\[
 E_{14}=E_4^2E_6,\qquad
 E_{16}=\frac{1617E_4^4+2000E_4E_6^2}{3617}
       =E_4\Delta\left(j-\frac{3456000}{3617}\right),
\]
which directly explain \eqref{newmodular:eq:F14}--\eqref{newmodular:eq:F16}.
The generator implements a standard modular-form construction, supplying
explicit reusable coefficients and proof certificates for the CM recipe.

\subsection{CM nonvanishing with a rational principal-lattice bound}
\label{cmnonzero:sec:nonvanishing}
The possible zero denominators in the product and genus formulas can be
excluded uniformly. Non-elliptic CM nonvanishing is classical, through
Kohnen's conjugation argument~\cite{cmnonzero:Kohnen}; see the introduction
of Chiriac--Jorza~\cite{cmnonzero:ChiriacJorza} for its place in zero theory.
The following proof supplies an explicit principal-point bound, certified
using rational arithmetic. It applies to nonmaximal orders as well.

\begin{theorem}[Nonvanishing at every non-elliptic CM point]
\label{cmnonzero:thm:allCM}
Let $D<-4$ be any quadratic-order discriminant. For every CM point $\tau$
of discriminant $D$ and every even integer $w\ge4$,
\[
 G_w(\tau)\ne0.
\]
For the principal reduced point
\[
 \tau_D=\begin{cases}i\sqrt{|D|}/2,&D\equiv0\pmod4,\\
 (-1+i\sqrt{|D|})/2,&D\equiv1\pmod4,\end{cases}
\]
one has the weight- and discriminant-uniform bound
$|G_w(\tau_D)|>7/20$.
\end{theorem}

\begin{proof}
The principal lattice has, for $n\ne0$, squared lengths bounded below by
one of the two integer quadratic forms
\[
 q_7(m,n)=m^2-mn+2n^2,\qquad q_8(m,n)=m^2+2n^2.
\]
Indeed the odd principal discriminant has $|D|\ge7$, and the even one
has $|D|\ge8$. Both forms are at least two when $n\ne0$. Hence for every
even $w\ge4$ the off-axis absolute sum is bounded by
\[
 T_d=\sum_{n\ne0}\sum_{m\in\mathbb Z}q_d(m,n)^{-2},\qquad d=7,8.
\]
We prove $T_7<33/20$ and $T_8<11/8$ by finite rational sums and explicit
infinite-tail bounds.

For any $A>0$ and real shift $s$, layer-cake integration gives
\[
 \sum_{m\in\mathbb Z}\frac1{((m-s)^2+A^2)^2}
 \le A^{-4}+\frac{\pi}{2A^3}.
\]
To see this, each superlevel set of the even decreasing function is an
interval. Its number of translated lattice points is at most its length
plus one. Integrating that inequality bounds the lattice sum by the
integral of the function plus its maximum. The integral is $\pi/(2A^3)$.

Use $N=10$, $M=20$ in the finite rectangle. The imaginary coordinate in
both lattices is greater than one. Thus the part with $|n|>N$, summed
over every $m$, is less than
\[
 \frac2{3N^3}+\frac2{N^2},
\]
by $\pi<4$, $\sum_{n>N}n^{-4}\le1/(3N^3)$ and
$\sum_{n>N}n^{-3}\le1/(2N^2)$. For $1\le |n|\le N$ and $|m|>M$,
the squared length is at least $(|m|-N/2)^2$. Integral comparison bounds
this part by $4N/[3(M-N/2)^3]$. The finite certificate therefore verifies
\[
 2\sum_{n=1}^{10}\sum_{m=-20}^{20}q_7(m,n)^{-2}
       +\frac{31}{1500}+\frac8{2025}<\frac{33}{20},
\]
and the same expression with $q_8$ is less than $11/8$. Every term and
comparison in these two finite inequalities is rational; the full fractions
are recorded in \texttt{verification/CM-nonvanishing.json}.

The $n=0$ axis contributes exactly $2\zeta(w)>2$. The triangle inequality
now gives $|G_w(\tau_D)|>2-33/20=7/20$ in the odd case, and the stronger
$5/8$ in the even case.

Suppose a CM value at the same discriminant vanished. By
\eqref{cmnorm:eq:Fw}, its singular modulus would be a zero of the rational
polynomial $F_w$. The class polynomial $H_D$ is the minimal polynomial of
that singular modulus, including for a nonmaximal order
\cite{cmnorm:Sutherland,cm:Cox}. Thus $H_D$ divides $F_w$, and the
principal singular modulus is also a zero. Since $\Delta$ never vanishes
on the upper half-plane, this would make $G_w(\tau_D)=0$, contradicting
the bound. Modular transformations multiply a positive-weight value by a
nonzero factor, so the conclusion covers every representative.
\end{proof}

Consequently all class and genus resultants in the following formulas are
nonzero at $D<-4$. Their positive roots and quotients are defined at every
even weight. The excluded discriminants $-3,-4$ retain their familiar
elliptic vanishing patterns.


\subsection{Class products in every even weight}
\begin{theorem}[CM product formula at every even weight]
\label{cmnorm:thm:allweights}
With the preceding definitions, for every even $w\geq4$,
\begin{align}
 \prod_{Q\in\mathcal Q_D}|G_w(\tau_Q)|
 &=\left(\frac{|B_w|}{w!}\right)^h
   \left(\frac{A_D}{d^h}\right)^{w/2}
   C_d^{w/4}(2\pi)^{w(2h-\varphi_d)/4}
   P_D^w\,|\mathcal R_{D,w}|^{1/12}.
 \label{cmnorm:eq:allweights}
\end{align}
Here the Bernoulli numbers satisfy $B_2=1/6$. Every factor and the resultant are nonzero by the preceding CM
nonvanishing theorem. The identity itself also applies to a zero product
when extended to an elliptic discriminant in the order formula below.
\end{theorem}

The absolute values are essential in a uniform statement. Selecting one
standard representative for each class need not make the product of
positive-weight modular forms real. Section~\ref{cmnorm:sec:phase} gives
the phase information needed for the manuscript's examples.

\subsection{Derivation of the period factor}

Put
\[
 R_D=\prod_{r=1}^{d-1}\Gamma(r/d)^{\chi_D(r)}.
\]
We first establish the precise normalization
\begin{equation}
 \prod_{Q\in\mathcal Q_D}|\Delta(\tau_Q)|
 =\frac{A_D^6 R_D^6}{(2\pi d)^{6h}}.
 \label{cmnorm:eq:delta}
\end{equation}
This is the fundamental-discriminant case of the classical
Lerch--Chowla--Selberg formula. For a modern primary treatment with
all normalization factors, see Barquero-Sanchez et al.\
\cite{cmnorm:Barquero} and Cohen\cite{cmnorm:Cohen2026}.
Here is a derivation that makes the powers of $2$, $\pi$, and $d$
checkable.

For $\operatorname{Re}s>1$, define the full Epstein sum
\[
 Z(\tau,s)=\sum_{(m,n)\neq(0,0)}
       \frac{(\operatorname{Im}\tau)^s}{|m+n\tau|^{2s}}.
\]
The class decomposition of the Dedekind zeta function gives
\begin{equation}
 \sum_{Q\in\mathcal Q_D} Z(\tau_Q,s)
 =2\left(\frac{\sqrt d}{2}\right)^s\zeta(s)L(s,\chi_D).
 \label{cmnorm:eq:epstein}
\end{equation}
Indeed $Q(m,-n)=a|m+n\tau_Q|^2$ (the sign change in $n$ does not
affect the full lattice sum), while
$a\operatorname{Im}\tau_Q=\sqrt d/2$; the factor $2$ is the
number of units, since $D<-4$.

The Kronecker limit formula at zero reads
\[
 Z(\tau,0)=-1,\qquad
 Z'(\tau,0)=-2\log\bigl(2\pi\sqrt{\operatorname{Im}\tau}
                               |\eta(\tau)|^2\bigr).
\]
Also $L(0,\chi_D)=h$, and the finite Hurwitz-zeta expression for $L$
and Lerch's identity give
\[
 L'(0,\chi_D)=-h\log d+\log R_D.
\]
Differentiate \eqref{cmnorm:eq:epstein} at zero, using
$\zeta(0)=-1/2$ and $\zeta'(0)=-\tfrac12\log(2\pi)$.
After collecting terms, one obtains
\begin{equation}
 \prod_Q\sqrt{\operatorname{Im}\tau_Q}|\eta(\tau_Q)|^2
 =(4\pi\sqrt d)^{-h/2}R_D^{1/2}.
 \label{cmnorm:eq:CS}
\end{equation}
Finally,
$\prod_Q\operatorname{Im}\tau_Q=(\sqrt d/2)^h/A_D$.
Dividing \eqref{cmnorm:eq:CS} by the square root of this product
and raising to the twelfth power proves \eqref{cmnorm:eq:delta}.

Because $\chi_D$ is odd, the positive and negative character residues
are interchanged by $r\mapsto d-r$. Euler reflection therefore gives
\begin{equation}
 R_D=\frac{P_D^2}{\pi^{\varphi_d/2}}
       \prod_{\chi_D(r)=1}\sin(\pi r/d)
     =\frac{P_D^2\sqrt{C_d}}{(2\pi)^{\varphi_d/2}}.
 \label{cmnorm:eq:reflection}
\end{equation}
To verify the last equality, square the sine product: it becomes the
product over all reduced residues, because paired sines are equal.
The identity $|1-e^{2\pi ir/d}|=2\sin(\pi r/d)$ and the cyclotomic
polynomial evaluated at $1$ give its square as $C_d/2^{\varphi_d}$.
Every sine is positive, fixing the square-root sign.

\begin{proof}[Proof of Theorem~\ref{cmnorm:thm:allweights}]
Multiply \eqref{cmnorm:eq:Fw} over the class representatives and take
absolute values:
\[
 \prod_Q|E_w(\tau_Q)|
 =\left(\prod_Q|\Delta(\tau_Q)|\right)^{w/12}
       |\mathcal R_{D,w}|^{1/12}.
\]
Substitute \eqref{cmnorm:eq:delta}, then
$2\zeta(w)=|B_w|(2\pi)^w/w!$ and
\eqref{cmnorm:eq:reflection}. Directly collecting exponents gives
\eqref{cmnorm:eq:allweights}. No logarithm is taken of $E_w$, so the
same argument applies when a factor vanishes.
\end{proof}

\subsection{A finite arithmetic recipe and useful low weights}

The general formula is particularly simple in the first few weights:
\begin{center}
\begin{tabular}{c l l}
\toprule
$w$ & $F_w(X)$ & $|\mathcal R_{D,w}|^{1/12}$\\
\midrule
$4$ & $X^4$ & $|H_D(0)|^{1/3}$\\
$6$ & $(X-1728)^6$ & $|H_D(1728)|^{1/2}$\\
$8$ & $X^8$ & $|H_D(0)|^{2/3}$\\
$10$ & $X^4(X-1728)^6$ & $|H_D(0)|^{1/3}|H_D(1728)|^{1/2}$\\
$12$ & $(X-432000/691)^{12}$ & $|H_D(432000/691)|$\\
\bottomrule
\end{tabular}
\end{center}
For the last row, use
$E_{12}=(441E_4^3+250E_6^2)/691
       =\Delta(j-432000/691)$.
This table is an exact algebraic calculation. In particular, the prime
factors of a class product are no longer guessed from a ratio: at weights
four and six they are determined by the factorizations of $H_D(0)$ and
$H_D(1728)$, together with the explicitly displayed elementary factors.

Specializing the theorem to weight four gives
\begin{equation}
 \prod_Q|G_4(\tau_Q)|
 =\frac{|H_D(0)|^{1/3}A_D^2C_d}
        {180^h d^{2h}2^{\varphi_d}\pi^{\varphi_d-2h}}P_D^4.
 \label{cmnorm:eq:weightfour}
\end{equation}
At weight six it gives
\begin{equation}
 \prod_Q|G_6(\tau_Q)|
 =\frac{A_D^3C_d^{3/2}|H_D(1728)|^{1/2}}
        {30240^h d^{3h}}
   (2\pi)^{3h-3\varphi_d/2}P_D^6.
 \label{cmnorm:eq:weightsix}
\end{equation}



\subsection{The phase and the exceptional class-pair average}
\label{cmnorm:sec:phase}

For a form strictly inside the reduced domain with $b\neq0$, the form
with $-b$ also occurs. Their $G_w$ values are complex conjugates, so their
product is positive unless a factor is zero. Boundary forms with $b=0$
or $b=a<c$ give real $G_4$ and $G_6$, and these values are positive.
For $G_4$ this follows from its positive value high on each boundary ray,
continuity, and the fact that $E_4$ has no zero there except the hexagonal
endpoint; for $G_6$ the same argument uses the unique elliptic zero at $i$.
The excluded discriminants $-3,-4$ are exactly the exceptional endpoints.
These are the usual zero-divisor consequences of the valence formula.

The remaining boundary type is $a=c$ and $0<b<a$. Then $|\tau_Q|=1$.
Modularity and conjugation imply
$\overline{E_4(\tau_Q)}=\tau_Q^4E_4(\tau_Q)$.
Following the unit arc from $i$ and using the absence of an interior zero
fixes the sign and yields
\begin{equation}
 G_4(\tau_Q)=-\tau_Q^{-2}|G_4(\tau_Q)|
 =\left(\frac{d-b^2}{4a^2}
         -\frac{ib\sqrt d}{2a^2}\right)|G_4(\tau_Q)|.
 \label{cmnorm:eq:arcphase}
\end{equation}
Thus an unqualified equality between a raw class product and its positive
gamma norm is not valid for arbitrary discriminants. Recent detailed work
on the analogous eta-phase issue is Cohen\cite{cmnorm:Cohen2026};
the elementary modular calculation above suffices for the present cases.

For $D=-15$ the nonprincipal form is $(2,1,2)$. Consequently
\[
 G_4((-1+i\sqrt{15})/4)
 =\frac{7-i\sqrt{15}}8
       |G_4((-1+i\sqrt{15})/4)|.
\]
The manuscript uses its real part, which is $7/8$ of the absolute value.
This explains its factor $7$ exactly; averaging a conjugate pair must
not be confused with taking one class representative in a norm.

\subsection{Proofs of the manuscript's five class products}

The needed class-polynomial constants and leading-coefficient products
are as follows. The supplied interval certificate proves these are the
exact Hilbert class polynomials, independently of numerical recognition.
\begin{center}
\begin{tabular}{c r r l l}
\toprule
$D$&$h$&$A_D$&$H_D(0)$&$H_D(1728)$\\
\midrule
$-15$&2&2&$-3^6 5^3 11^3$&$3^6 7^4 11^2$\\
$-20$&2&2&$-2^{12}5^3 11^3$&$-2^{16}11^2 19^2$\\
$-23$&3&4&$5^9 11^3 17^3$&$7^6 11^4 19^2 23$\\
$-39$&4&12&$3^{15}17^3 23^3 29^3$&$3^{12}7^8 19^4 23^2$\\
$-47$&5&36&$5^{15}11^6 23^3 29^3$&$11^6 19^4 23^4 31^2 43^2 47$\\
\bottomrule
\end{tabular}
\end{center}
The signs and factorizations in this table are exact integer arithmetic;
they are also recorded in the machine receipt.

\begin{corollary}[The five recorded class products are proved]
\label{cmnorm:cor:recorded}
The product formulas in Theorems~\ref{cm:res:d15}, \ref{cm:res:d20},
\ref{cm:res:d23}, \ref{cm:res:d39} and \ref{cm:res:d47} hold exactly,
with the real-part projection specified at discriminant $-15$.
\end{corollary}
\begin{proof}
For $D=-23,-39,-47$ there are no unit-arc boundary representatives.
Interior pairs are conjugates and all remaining weight-four boundary
values are positive, so the absolute norm is the raw product. Substitute
the factorizations in the preceding table into
\eqref{cmnorm:eq:weightfour}; these give exactly the three displayed
class products. For $D=-15$, the real-part factor $7/8$ in
\eqref{cmnorm:eq:arcphase} gives the recorded factor $7$.

For $D=-20$, the norm first gives
\[
 \prod_QG_4(\tau_Q)=\frac{11P_{-20}^4}{2^{14}3^4 5^5\pi^4}.
\]
The multiplication and reflection formulas give
\[
 \frac{\Gamma(1/20)\Gamma(9/20)}{\Gamma(3/20)\Gamma(7/20)}
 =4\,5^{1/4}\sin(3\pi/20)\sin(7\pi/20)=5^{1/4}\varphi,
\]
where $\varphi=(1+\sqrt5)/2$. Indeed apply the fivefold multiplication
formula at $1/20$ and reflect the factors at $13/20$ and $17/20$.
Thus $P_{-20}=5^{1/4}\varphi\Gamma(3/20)^2\Gamma(7/20)^2$;
substitution proves the remaining displayed formula. No period
independence or numerical recognition is used in this deduction.
\end{proof}

\subsection{Three explicit weight-six extensions}

The same exact class data now produce formulas that were not displayed
in the baseline manuscript:
\begin{align}
 \prod_{Q\in\mathcal Q_{-23}}G_6(\tau_Q)
 &=\frac{11^2\cdot19\,P_{-23}^6}
       {2^{33}3^9 5^3 23^7\pi^{24}},\label{cmnorm:eq:G6d23}\\
 \prod_{Q\in\mathcal Q_{-39}}G_6(\tau_Q)
 &=\frac{19^2\cdot23\,P_{-39}^6}
       {2^{38}3^{15}5^4 13^{12}\pi^{24}},\label{cmnorm:eq:G6d39}\\
 \prod_{Q\in\mathcal Q_{-47}}G_6(\tau_Q)
 &=\frac{11^3 19^2 23^2\cdot31\cdot43\,P_{-47}^6}
       {2^{73}3^9 5^5 7^5 47^{13}\pi^{54}}.
       \label{cmnorm:eq:G6d47}
\end{align}
They follow by inserting the factorizations of $H_D(1728)$ in
\eqref{cmnorm:eq:weightsix}; positivity was already established.
For example, at $D=-23$ the square root is
$7^3 11^2 19\sqrt{23}$, and $C_{23}^{3/2}=23^{3/2}$.
Their product contributes $23^2$, which changes $23^9$ in the
denominator to $23^7$. This calculation is a useful check on the
normalization.

Every factor here is a polygamma row sum. For example,
\[
 G_6(\tau)=2\zeta(6)+\frac1{60}
  \sum_{n\geq1}\left[\psi^{(5)}(n\tau)
                         +\psi^{(5)}(1-n\tau)\right].
\]
Consequently \eqref{cmnorm:eq:G6d23}--\eqref{cmnorm:eq:G6d47}
are exact identities for finite products of exponentially convergent
polygamma-row combinations. Weight twelve gives a further family whose
entire arithmetic factor is the rational number
$|H_D(432000/691)|$.

\subsection{How the class-polynomial data are certified}
\label{cmnorm:sec:certification}

The program \texttt{cm\_norms.py} enumerates reduced forms by exact
integer arithmetic. It then evaluates all singular moduli in rectangular
complex intervals using 120 decimal digits of directed interval
arithmetic and 70 terms of the Lambert series
\[
 E_4=1+240\sum_{n\geq1}\frac{n^3q^n}{1-q^n},\qquad
 E_6=1-504\sum_{n\geq1}\frac{n^5q^n}{1-q^n},\qquad
 j=\frac{1728E_4^3}{E_4^3-E_6^2}.
\]
The interval engine encloses $\pi$, $\sqrt d$, and the complex
exponential. Its output endpoints are finite binary rationals; the
program converts them to exact Python fractions before every final
containment comparison.

Reduced points satisfy $\operatorname{Im}\tau\geq\sqrt3/2$, hence
$|q|\leq e^{-\pi\sqrt3}<1/200$. This last strict inequality needs
no floating-point assumption: use $\pi>314/100$,
$\sqrt3>173/100$, their product greater than $27/5$, and the rational
partial sum $\sum_{k=0}^{15}(27/5)^k/k!>200$.
For $r=1/200$ and $p=3,5$, the omitted tail satisfies
\begin{equation}
 \left|\sum_{n>N}\frac{n^pq^n}{1-q^n}\right|
 \leq\frac{(N+1)^p r^{N+1}\mathcal A_p(r)}{(1-r)^{p+2}},
 \label{cmnorm:eq:tail}
\end{equation}
where
\[
 \mathcal A_3(r)=1+4r+r^2,\qquad
 \mathcal A_5(r)=1+26r+66r^2+26r^3+r^4.
\]
Indeed $n=N+1+j\leq(N+1)(j+1)$, and
$\sum_{j\geq0}(j+1)^pr^j=\mathcal A_p(r)/(1-r)^{p+1}$.
The code adds the resulting error bound in both real and imaginary
coordinates before forming $j$ or multiplying the class polynomial.

Classical CM theory proves in advance that each coefficient of $H_D$
is an integer. The interval product encloses every coefficient in an
interval strictly contained in its asserted integer plus
$(-1/2,1/2)$; its imaginary interval contains zero. Thus the unique
possible integer is proved, rather than merely recognized to many digits.
The maximum real interval widths were
\begin{center}
\begin{tabular}{c c c c c}
\toprule
$D=-15$&$D=-20$&$D=-23$&$D=-39$&$D=-47$\\
\midrule
$1.36\cdot10^{-109}$&$3.95\cdot10^{-108}$&$1.83\cdot10^{-103}$&
$2.80\cdot10^{-95}$&$1.51\cdot10^{-88}$\\
\bottomrule
\end{tabular}
\end{center}
The displayed bounds have been rounded upward. Exact interval endpoints,
all polynomial coefficients, and the integer factorizations are saved
in \texttt{cm\_norms\_receipt.json}.

An additional 25 numerical comparisons test the gamma norm formula at
weights $4,6,8,10,12$ for each of the five discriminants. They have
relative residuals below $1.1\cdot10^{-119}$ at 120 digits. These
comparisons are regression checks: they are not interval proofs of the
gamma identities, whose proof is Theorem~\ref{cmnorm:thm:allweights}.
Similarly, the certified-polynomial computation is a numerical proof
conditional on the standard CM integrality theorem and on the interval
library's directed-rounding implementation; it is not a Lean
formalization. This evidence boundary is explicit and reproducible.

\subsection{Extension to arbitrary quadratic orders}

There is also a direct version for $D=D_0f^2$, where $D_0$ is fundamental,
including the unit exceptions $D_0=-3,-4$. Let $h=h(D)$ and
$d=|D|$, and let $\mathrm{CS}(D_0)$ be Cohen's normalized gamma
quotient
\[
 \mathrm{CS}(D_0)=
 \left(\prod_{r=1}^{|D_0|-1}\Gamma(r/|D_0|)^{\chi_{D_0}(r)}
 \right)^{w(D_0)/(2h(D_0))}.
\]
Define the explicit conductor factor
\[
 c(D_0,f)=\prod_{p\mid f}p^{e(p)},\qquad
 e(p)=\frac{(1-p^{-v_p(f)})(1-\chi_{D_0}(p))}
 {(1-1/p)(p-\chi_{D_0}(p))}.
\]
The order form of Lerch--Chowla--Selberg\cite{cmnorm:Barquero,cmnorm:Cohen2026}
gives the following immediate extension of the proof:
\begin{equation}
 \prod_{Q\in\mathcal Q_D}|G_w(\tau_Q)|
 =(2\zeta(w))^h A_D^{w/2}
 \left(\frac{c(D_0,f)\mathrm{CS}(D_0)}{2\pi d}\right)^{hw/2}
 |\operatorname{Res}(H_D,F_w)|^{1/12}.
 \label{cmnorm:eq:orders}
\end{equation}
Thus nonmaximal orders require an explicit conductor correction, not a
new unidentified period basket. Setting $f=1$ and $D<-4$ recovers the
main theorem.

\begin{corollary}[The four class-number-one row evaluations]
\label{cmnorm:cor:classone}
The four formulas at $2i$ and $i\sqrt2$ in
Theorem~\ref{cm:res:2i} hold exactly.
\end{corollary}
\begin{proof}
The additional coefficient certificates give
$H_{-8}(X)=X-8000$ and $H_{-16}(X)=X-287496$, each with one reduced
class. At $D=-8$ insert $H_{-8}(0)$ and $H_{-8}(1728)$ into the
fundamental formula; $C_8=2$ and
$P_{-8}=\Gamma(1/8)\Gamma(3/8)$. This gives the two recorded coefficients.
At $D=-16$ use \eqref{cmnorm:eq:orders} with $D_0=-4$, $f=2$,
$c(-4,2)=\sqrt2$ and
\[
 \mathrm{CS}(-4)=\left(\frac{\Gamma(1/4)}{\Gamma(3/4)}\right)^2
              =\frac{\Gamma(1/4)^4}{2\pi^2}.
\]
The weight-four and weight-six resultant roots are
$287496^{1/3}=66$ and
$(287496-1728)^{1/2}=378\sqrt2$. Substitution gives the stated factors
$11/15360$ and $1/(2^{14}5)$. These principal imaginary-axis row sums
are positive by the principal-point proof, so absolute values may be removed.
The two new class polynomials are certified by directed intervals and CM
integrality at 120 digits and 70 Lambert terms, with exact recorded endpoints.
\end{proof}

\begin{proposition}[Two exact Weber degree statements]
\label{cmweber:prop:degrees}
At discriminant $-23$, the phase-corrected cubic coordinate defined in
Theorem~\ref{cm:res:d23poly} is unique and belongs to the same cubic field
as its singular modulus. At discriminant $-39$, every cube root of a
singular modulus has degree twelve over $\Q$, and none belongs to the
Hilbert class field of degree eight over $\Q$.
\end{proposition}
\begin{proof}
Exact polynomial arithmetic gives
\[
 \operatorname{Res}_Y(Y^3+155Y^2+650Y+23375,\,Y^3-X)=-H_{-23}(X).
\]
The three cube images are therefore the three distinct singular moduli.
For each one exactly one root of the cubic has that image. Since
$[\Q(j_Q):\Q]=3$ and $j_Q=\xi_Q^3$, the cubic equation forces
$\Q(\xi_Q)=\Q(j_Q)$. The quotient of $\xi_Q$ by $E_4/\eta^8$ is
necessarily a cube root of unity, giving the precise phase correction.

For $D=-39$, the integer polynomial $H_{-39}(Y^3)$ reduces modulo five to
\[
 f(Y)=Y^{12}+Y^9+4Y^6+4Y^3+3.
\]
The independent finite-field certificate verifies
\[
 Y^{5^{12}}\equiv Y\pmod f,\qquad
 \gcd(f,Y^{5^6}-Y)=\gcd(f,Y^{5^4}-Y)=1.
\]
Rabin's criterion proves $f$ irreducible over $\mathbf F_5$; its monic
integer lift is thus irreducible over $\Q$. Every possible phase has
cube $j_Q$, so it is a root of this degree-twelve polynomial. Such a
field cannot be a subfield of the degree-eight Hilbert class field.
The resultant, the modular witness and its independent standard-library
Frobenius/gcd replay are retained in the CM verification receipts.
\end{proof}

\subsection{Research directions and scope}

The theorem suggests several concrete continuations.
\begin{enumerate}
\item Extend the exact polynomial $F_w$ generator supplied here and combine it with
certified class-polynomial software such as Sutherland's CRT methods
\cite{cmnorm:Sutherland}. This extends the product database to large
weights and class numbers without integer-relation searches.
\item Determine signed, rather than absolute, products for all weights
by tracking the phases of boundary representatives and the zero signs
of $E_w$ along the unit arc. The present phase proof covers the displayed
weight-four and weight-six examples; a general norm alone does not
settle this problem.
\item Use factorizations of $H_D(0)$ and $H_D(1728)$ to compare the prime
support with arithmetic formulas for norms of singular-modulus
differences. This gives an explanation for the prime patterns observed
in the manuscript, beyond using those patterns to build a PSLQ basis.
\item Extend the certified computations to class-by-class algebraic
factors and genus-wise products. A full class norm does not by itself
prove class ratios. Section~\ref{cmgenus:sec:main} supplies the additional
genus-character argument for the recorded ratios at $D=-15,-20,-39$;
extending it beyond these examples remains a natural direction.
\item Build a uniform certificate format containing the discriminant,
reduced forms, exact $H_D$, exact $F_w$, resultant, conductor factor,
phase convention, gamma residue set, and analytic tail bound. Each
claimed identity then has a finite, reviewable provenance chain.
\end{enumerate}

None of these norm computations proves linear or algebraic independence
of individual gamma values or individual polygamma rows. Nor does it
imply that arbitrary sums of raw positive-weight CM values are class
invariants. The result is a precise product theorem with completely
specified representative dependence.
